\documentclass[12pt,a4paper]{article}
\usepackage{cmap} % Makes the PDF copiable. See http://tex.stackexchange.com/a/64198/25761
\usepackage[T1]{fontenc}
\usepackage[brazil]{babel}
\usepackage[utf8]{inputenc}
\usepackage{amsmath}
\usepackage{amsfonts}
\usepackage{amssymb}
\usepackage{amsthm}
\usepackage{textcomp} % \degree
\usepackage{gensymb} % \degree
\usepackage[usenames,svgnames,dvipsnames]{xcolor}
\usepackage{hyperref}
\usepackage{graphicx}
\usepackage[margin=2cm]{geometry}

\hypersetup{
    colorlinks = true,
    allcolors = {blue}
}

% TODO: Consider using exsheets
% http://linorg.usp.br/CTAN/macros/latex/contrib/exsheets/exsheets_en.pdf
%
% http://ctan.org/tex-archive/macros/latex/contrib/exercise/
% Options: answerdelayed,lastexercise,noanswer
\usepackage[answerdelayed,lastexercise]{exercise}

\addto\captionsbrazil{%
\def\listexercisename{Lista de exerc\'icios}%
\def\ExerciseName{Exerc\'icio}%
\def\AnswerName{Solu\c{c}\~ao do exerc\'icio}%
\def\ExerciseListName{Ex.}%
\def\AnswerListName{Solu\c{c}\~ao}%
\def\ExePartName{Parte}%
\def\ArticleOf{de\ }%
}

\renewcommand{\ExerciseHeaderTitle}{(\ExerciseTitle)\ }
\renewcommand{\ExerciseListHeader}{%\ExerciseHeaderDifficulty%
\textbf{%\ExerciseListName\
\ExerciseHeaderNB.\ %
%\ --- \
\ExerciseHeaderTitle}%
%\ExerciseHeaderOrigin
\ignorespaces}
\renewcommand{\AnswerListHeader}{\textbf{\ExerciseHeaderNB.\ (\AnswerListName)\ }}

\newcommand{\fixme}{{\color{red}(...)}}

\newcommand*\sen{\operatorname{sen}}
\newcommand*\tg{\operatorname{tg}}
\newcommand*\R{\mathbb{R}}

\renewcommand{\theenumi}{\alph{enumi}}
\renewcommand\labelenumi{(\theenumi) }

\newcommand*\tipo{Prova IV}
\newcommand*\turma{TADS121-01C}
\newcommand*\disciplina{CDI0001}
\newcommand*\eu{Helder G. G. de Lima}
\newcommand*\data{12/06/2017}

\author{\eu}
\title{\tipo - \disciplina}
\date{\data}

\begin{document}
\thispagestyle{empty}
\newgeometry{margin=2cm,bottom=0.5cm}
\begin{center}
\includegraphics[width=9.0cm]{marca} \\
\textbf{\tipo\ (\disciplina / \turma)} \\
Prof. \eu\footnote{
Este é um material de acesso livre distribuído sob os termos da licença \href{https://creativecommons.org/licenses/by-sa/4.0/deed.pt_BR}{Creative Commons Atribuição-CompartilhaIgual 4.0 Internacional}}
\end{center}

\noindent Nome do(a) aluno(a): \underline{\hspace{9,7cm}} Data: \underline{\data}

\begin{center}\fbox{
\begin{minipage}{14cm}

{\footnotesize
\begin{itemize}
\renewcommand{\theenumi}{\Roman{enumi}}
\item Identifique-se em todas as folhas.
\item Mantenha o celular e os demais equipamentos eletrônicos desligados durante a prova.
\item Justifique cada resposta com cálculos ou argumentos baseados na teoria estudada.
\item Resolva (integralmente) apenas as questões de que precisar para somar 10,0 pontos.
\end{itemize}
}

\end{minipage}
}
\end{center}

\begin{ExerciseList}
\Exercise[title={2,0}] Seja $y = f(x)$ uma função derivável definida implicitamente por $\tg(\pi y) + x(x - 2) = 0$. Obtenha a(s) reta(s) tangente(s) ao gráfico de $f$ no(s) ponto(s) em que a ordenada $y = 0$.
\Answer \fixme

\Exercise[title={2,0}] Um avião passou por você a uma altura de $10 km$ e uma velocidade constante de $910\ km/h$. Quão rápido a distância entre você e o avião estará aumentando quando ele passar sobre um parque que fica exatamente a $24 km$ de onde você o observa? Assuma que o avião se desloca em linha reta, a uma altura constante e com uma velocidade constante.
\Answer \fixme

\Exercise[title={2,0}] Qual é o mínimo de material necessário para confeccionar uma caixa em forma de paralelepípedo com base quadrada, cujo volume seja de $8$ litros ($8000 cm^3$)?
\Answer \fixme

\Exercise[title={2,0}] Analize se $f(x) = \dfrac{\ln(x^4) + x^3 + 4}{x^2}$ possui assíntotas oblíquas e se possui assíntotas verticais. Quando for o caso, exiba as equações de tais assíntotas.

\Exercise[title={4,0}] Seja $f(x) = \dfrac{3 - x^2}{3 + x^2}$. Utilize as ferramentas do cálculo diferencial para:
\begin{enumerate}
\item (0,4) Obter as derivadas de primeira e segunda ordem da função $f$.
\item (0,6) Obter $\lim_{x \to \pm\infty} f(x)$.
\item (1,0) Obter todos os intervalos $(a,b)$ em que a função:
\begin{enumerate}
\item É \textbf{crescente} e \textbf{decrescente}, respectivamente.
\item Tem \textbf{concavidade para cima} e \textbf{para baixo}, respectivamente.
\end{enumerate}
\item (1,0) Identificar no domínio da função, se existirem, todos os pontos \textbf{críticos}, de \textbf{máximo} e de \textbf{mínimo} (locais e globais) e de \textbf{inflexão}. Justifique suas conclusões com base nas informações que obteve nos itens (a), (b) e (c).
\item (1,0) Esboçar o \textbf{gráfico} da função, usando (corretamente) os dados anteriores.
\end{enumerate}
\Answer \fixme
\end{ExerciseList}

\begin{center}
BOA PROVA!
\end{center}

%\newpage
%\restoregeometry
%\section*{Respostas}
%\shipoutAnswer
\end{document}
